\documentclass[11pt,a4paper]{article}
\usepackage[margin=1in]{geometry}
\usepackage{amsmath,amssymb,amsthm,mathtools}
\usepackage{booktabs}
\usepackage[hidelinks]{hyperref}

\newtheorem{definition}{Definition}
\newtheorem{lemma}{Lemma}
\newtheorem{proposition}{Proposition}
\newtheorem{theorem}{Theorem}
\newtheorem{corollary}{Corollary}
\newtheorem{remark}{Remark}

\title{\textbf{Static Ledger: Boolean Inverse Algebra over \(GF(4)\)}\\[0.4em]
\large Exact Complement, Gate Reduction, and Inverse-Metric Closure}
\author{Foliation Engine --- Theoretical Formalization (W10)}
\date{July 2026}

\begin{document}
\maketitle

\begin{abstract}
We record a purely algebraic settlement of finite-field Boolean inversion on
\(GF(4)\cong GF(2^2)\). Sequence symbols are identified with field elements;
the exact complement equation \(\mathbf{S}_{\mathrm{path}}\oplus\mathbf{S}_{\mathrm{thera}}=\mathbf{1}\)
is solved by a depth-\(O(1)\) Boolean map on the natural \(GF(2)^2\) bit lift;
macrocyclic closure is stated as an exact residual \(\Delta E=0\) under an
inverse Hermitian metric constraint. No wet-lab protocol, no statistical
estimator, and no affinity model appear in the ledger.
\end{abstract}

\section{Galois Field and Symbol Isomorphism}

\begin{definition}[Binary extension field]
Let \(\mathbb{F}_2=\{0,1\}\) and let
\[
  GF(4)=GF(2^2)=\mathbb{F}_2[x]/(x^2+x+1).
\]
Write \(\alpha\) for the residue class of \(x\). Then
\(\{0,1,\alpha,\alpha^2\}\) is the additive group of \(GF(4)\), with
\(\alpha^2+\alpha+1=0\) and \(\alpha^3=1\).
\end{definition}

\begin{definition}[Canonical symbol isomorphism]
Define \(\iota:\{A,U,C,G\}\to GF(4)\) by
\begin{equation}
  \iota(A)=0,\qquad
  \iota(U)=1,\qquad
  \iota(C)=\alpha,\qquad
  \iota(G)=\alpha^2.
\end{equation}
Equivalently, under the ordered bit lift
\(\varphi:GF(4)\to GF(2)^2\) given by
\(\varphi(b_0+b_1\alpha)=(b_0,b_1)\),
\begin{equation}
\begin{aligned}
  \iota(A)&\mapsto(0,0),&
  \iota(U)&\mapsto(1,0),\\
  \iota(C)&\mapsto(0,1),&
  \iota(G)&\mapsto(1,1).
\end{aligned}
\end{equation}
\end{definition}

\begin{remark}
The map \(\iota\) is a set bijection, not a ring homomorphism from a
nucleotide monoid. All subsequent identities are identities in
\(GF(4)^N\) or in the bit lift \(GF(2)^{2N}\).
\end{remark}

\begin{definition}[Path vector]
For length \(N\), a path state is a vector
\(\mathbf{S}_{\mathrm{path}}=(s_i)_{i=0}^{N-1}\in GF(4)^N\).
Its bit lift is
\(\mathbf{B}=(b_{i,0},b_{i,1})_{i=0}^{N-1}\in GF(2)^{2N}\) with
\(\varphi(s_i)=(b_{i,0},b_{i,1})\).
\end{definition}

\section{Exact-Form Boolean Complement}

\begin{definition}[Field XOR and the ledger one-vector]
Write \(\oplus\) for coordinatewise addition in \(GF(4)\).
Let the \emph{ledger one-vector} be the constant vector of the primitive
element \(\alpha\):
\begin{equation}
  \mathbf{1}:=(\alpha,\alpha,\ldots,\alpha)^\top\in GF(4)^N.
\end{equation}
(The multiplicative unit of \(GF(4)\) is written \(1_{\times}\) when needed;
it is \emph{not} the target of~\eqref{eq:complement}.)
\end{definition}

\begin{theorem}[Unique Boolean inverse]
For every \(\mathbf{S}_{\mathrm{path}}\in GF(4)^N\) there exists a unique
\(\mathbf{S}_{\mathrm{thera}}\in GF(4)^N\) such that
\begin{equation}
  \label{eq:complement}
  \mathbf{S}_{\mathrm{path}}\oplus\mathbf{S}_{\mathrm{thera}}=\mathbf{1}.
\end{equation}
Explicitly,
\[
  \mathbf{S}_{\mathrm{thera}}=\mathbf{1}\oplus\mathbf{S}_{\mathrm{path}}.
\]
\end{theorem}

\begin{proof}
\((GF(4),+)\) is an elementary abelian \(2\)-group, so each coordinate
equation \(s_i\oplus t_i=\alpha\) has the unique solution
\(t_i=\alpha\oplus s_i\). Vector uniqueness follows coordinatewise.
\end{proof}

\begin{proposition}[Bit-level gate reduction]
Under the bit lift \(\varphi(b_0+b_1\alpha)=(b_0,b_1)\),
equation~\eqref{eq:complement} is equivalent to the Boolean system
\begin{equation}
  \label{eq:gates}
  t_{i,0}=b_{i,0},\qquad
  t_{i,1}=\neg b_{i,1}
  \qquad(i=0,\ldots,N-1),
\end{equation}
where \(\neg\) denotes complement in \(GF(2)\) and
\(\varphi(t_i)=(t_{i,0},t_{i,1})\).
\end{proposition}

\begin{proof}
Write \(s_i=b_{i,0}+b_{i,1}\alpha\). Then
\[
  \alpha\oplus s_i
  =(0\oplus b_{i,0})+(1\oplus b_{i,1})\alpha
  =b_{i,0}+(\neg b_{i,1})\,\alpha,
\]
hence \(\varphi(\alpha\oplus s_i)=(b_{i,0},\neg b_{i,1})\), which is
exactly~\eqref{eq:gates}. The four-row truth table is:
\[
\begin{array}{c|cc|c}
s & (b_0,b_1) & \alpha\oplus s & (t_0,t_1)\\
\hline
0 & (0,0) & \alpha & (0,1)\\
1_{\times} & (1,0) & \alpha^2 & (1,1)\\
\alpha & (0,1) & 0 & (0,0)\\
\alpha^2 & (1,1) & 1_{\times} & (1,0).
\end{array}
\]
\end{proof}

\begin{corollary}[Combinational complexity]
The map \(\mathbf{S}_{\mathrm{path}}\mapsto\mathbf{S}_{\mathrm{thera}}\) is a
coordinatewise Boolean function of circuit depth \(O(1)\) and size \(O(N)\).
No iteration or search is required.
\end{corollary}

\section{Inverse-Metric Constraint and Exact Closure}

\begin{definition}[Clearance slack and inverse metric]
Let \(q\in\mathbb{R}^{3m}\) collect Cartesian coordinates of \(m\) atoms in a
closed macromolecular graph. For each steric pair \(e\) with clearance
slack \(\sigma_e(q)>0\), define the local metric weight
\begin{equation}
  g_e(q)=\frac{\mu L^2}{\sigma_e(q)^2},\qquad
  \mu>0,\ L>0,
\end{equation}
and the inverse weight \(g_e^{-1}(q)=1/g_e(q)\). Collecting weights into a
diagonal metric \(G(q)\) yields the inverse
\begin{equation}
  G^{-1}(q)=\mathrm{diag}\bigl(g_e^{-1}(q)\bigr).
\end{equation}
This \(G^{-1}\) is the static ledger's inverse Hermitian (K\"ahler-type)
brake: \(g_e^{-1}\to 0\) as \(\sigma_e\to 0\).
\end{definition}

\begin{definition}[Closure residual]
Let \(\mathcal{T}\) embed a \(GF(4)\) state into a torsional chart of the
closed graph, and let \(E(q)\) denote the exact-form cycle residual
(discrete circulation of edge lengths about each independent cycle).
Exact closure is the null statement
\begin{equation}
  \label{eq:deltaE}
  \Delta E := E(q)=0.
\end{equation}
\end{definition}

\begin{theorem}[Coupled algebraic--geometric settlement]
Suppose \(\mathbf{S}_{\mathrm{thera}}\) satisfies~\eqref{eq:complement} and
\(q^\star\) lies in the feasible open set \(\{\sigma_e>0\}\). If the inverse
metric annihilates the coupled tensor chart,
\begin{equation}
  \label{eq:kahler}
  \mathcal{H}^{-1}\!\left(
    \mathcal{T}(\mathbf{S}_{\mathrm{path}})
    \otimes
    \mathcal{T}(\mathbf{S}_{\mathrm{thera}})
  \right)
  =\mathbf{0},
\end{equation}
and the discrete cycle operator enforces~\eqref{eq:deltaE}, then the pair
\((\mathbf{S}_{\mathrm{thera}},q^\star)\) is an exact algebraic--geometric
settlement: Boolean complement holds identically, and macrocyclic closure
holds with residual zero.
\end{theorem}

\begin{proof}
Equation~\eqref{eq:complement} is already settled by Theorem~1.
Feasibility \(\sigma_e(q^\star)>0\) keeps \(G^{-1}\) finite. Condition
\eqref{eq:kahler} forces the coupled chart into the kernel of the inverse
metric brake. Independently, \(\Delta E=0\) is the definition of exact-form
null for cyclic edge residuals. The conjunction is therefore a static
settlement certificate in the product space
\(GF(4)^N\times\mathbb{R}^{3m}\).
\end{proof}

\begin{proposition}[No affinity content]
The identity \(\Delta E=0\) is a geometric closure constraint. It does not
imply a free-energy ranking, a scoring-function minimum, or a statistical
association constant.
\end{proposition}

\section{Static Settlement Ledger}

\begin{center}
\begin{tabular}{@{}lll@{}}
\toprule
Object & Symbol & Status\\
\midrule
Field & \(GF(4)\cong\mathbb{F}_2[x]/(x^2+x+1)\) & fixed\\
Isomorphism & \(\iota\) & fixed\\
Complement & \(\mathbf{S}_{\mathrm{path}}\oplus\mathbf{S}_{\mathrm{thera}}=\mathbf{1}\) & unique\\
Gate reduction & \(t_{i,0}=b_{i,0},\ t_{i,1}=\neg b_{i,1}\) & \(O(1)\) depth\\
Inverse metric & \(G^{-1}\to 0\) as slack \(\to 0\) & brake\\
Closure & \(\Delta E=0\) & exact null\\
\bottomrule
\end{tabular}
\end{center}

\section*{Disclaimer}
This document is a mathematical ledger only
(\texttt{clinical\_grade=false}). It encodes finite-field identities and
geometric closure constraints. It does not claim therapeutic effect,
manufacturing readiness, or empirical binding.

\end{document}
